\section{Problem Statement}
\label{sec:problem}

Developers reuse agent skills by copying a skill's folder into their own repository, often
editing it on the way in. No registry entry, manifest, or version pin records this
relationship, so once a skill is copied its lineage is lost. Three groups pay for this.
Developers who adopt a skill are never told when its source is fixed or improved, and because
agents execute skills, a stale copy keeps acting on outdated or unsafe instructions. Skill
authors cannot see where their work has spread or how it was changed. Security reviewers who
flag one harmful skill cannot locate its edited relatives.

Existing studies address this only in part. Approaches that trace reuse need an explicit
record of it: either the copy still carries the name of a skill published on a public
registry~\cite{gao2026registry}, or it lives in a repository created with GitHub's fork
feature~\cite{wu2026adaptation}. Copies that were renamed, heavily edited, or taken from an
unregistered repository are invisible to both. Clone detection for
skills~\cite{zhu2026skillreuse} removes the need for such records, but it measures
similarity, not origin. It cannot tell a derived variant from a template sibling, it does not
order related skills in time, and it has not been evaluated on the highly similar but
independently written pairs that templates produce.

The result is a gap in what we know. The only study that measured whether copies receive
updates found that 40.2\% of never-edited reused skills had missed later changes to their
source~\cite{gao2026registry}, but that measurement covers only the registry-linked minority
of skills. For skills shared by plain copying, which GitSkills shows to be
widespread~\cite{destefanis2027gitskills}, we cannot yet say which skills derive from which,
how concentrated reuse really is once edited copies are counted, or whether changes to an
original ever reach its copies.
